This study proposes a lightweight lithium-ion battery SOH estimation framework for resource-constrained battery management systems to address limited cross-cell HI stability and the difficulty of balancing prediction accuracy with computational efficiency. The study focuses on systematic HI construction and lightweight network design. The proposed multi-source health indicator extraction and optimization algorithm first constructs multiple types of candidate HIs from charging and discharging data and their derived curves. Based on correlations on the feature-development cell set, MS-CCCT then adaptively calibrates the constant-current charging voltage window at multiple scales, and PCC/SCC dual-threshold admission and redundancy removal determine the model inputs. With feature definitions and parameters held fixed, the selected HIs retain strong linear and monotonic relationships with SOH on other cells within the same dataset. For sequence modeling, MS-AgentNet combines ReLU² agent attention with small- and large-kernel depthwise separable convolutions to enhance the representation of local degradation variations and degradation trends over longer time scales through local feature extraction and global context modeling. With a fixed number of agents, RAA uses a small number of static learnable agents for information aggregation and broadcasting, reducing the theoretical complexity of attention-based correlation interactions from $O(N^2d)$ to $O(Nn_a d)$.

Experiments on four public battery datasets, Oxford, CALCE CS2, CALCE CX2, and MIT/Severson, show that MS-AgentNet achieves better overall performance than the four baseline models, CNN-Transformer, CNN-LSTM, Transformer, and LSTM. Specifically, the average MAPE of MS-AgentNet over the two cells in the CALCE CX2 dataset is 52.69\% lower than that of CNN-Transformer. Under the standardized complexity-test settings, MS-AgentNet also has low forward-pass computation, a small parameter count, and a small weight storage footprint. Ablation studies further show complementary effects between multi-scale DSConv and RAA and between DSConv-S and DSConv-L, contributing to more consistent overall performance across degradation scenarios. Overall, the proposed framework achieves a good balance between SOH estimation accuracy and resource consumption. It maintains good estimation performance on the test cells in the corresponding datasets, showing a degree of in-domain cross-cell generalization and potential for application in resource-constrained battery management systems.

Despite these results, the framework still relies on identifiable charging and discharging data segments, and its cross-dataset adaptation performance is affected by domain differences, source-domain selection, and the amount of target-domain data. Future work will focus on HI construction for incomplete charging and dynamic operating segments, together with lightweight domain adaptation strategies. Further validation will cover more battery chemistries, real-vehicle operating data, and embedded hardware platforms to improve the framework's practical applicability and deployment reliability.
